\documentclass[10pt,a4paper]{article}
\usepackage[margin=1.8cm]{geometry}
\usepackage{amsmath,amssymb}
\usepackage{enumitem}
\usepackage{booktabs}
\usepackage{xcolor}
\usepackage{titlesec}
\usepackage{hyperref}
\usepackage{multicol}

\definecolor{teal}{RGB}{0,128,128}
\definecolor{lightgray}{RGB}{245,245,245}

\titleformat{\section}{\large\bfseries\color{teal}}{}{0em}{}[\titlerule]
\titleformat{\subsection}{\normalsize\bfseries}{}{0em}{}

\setlength{\parindent}{0pt}
\setlength{\parskip}{4pt}

\title{\textbf{Q\&A Cheatsheet --- Adaptive Algorithms}\\[4pt]
\large Continuous Optimization $\cdot$ Project Presentation $\cdot$ 20.05.2026}
\author{Erion Krasniqi \& Denis Mustafa Xhabrahimi}
\date{}

\begin{document}
\maketitle
\vspace{-1em}

%=============================================================================
\section{Mathematical Definitions}
%=============================================================================

\subsection{Q: ``What does it mean for a function to be convex?''}
A function $f$ is convex if for all $x, y$ and $\lambda \in [0,1]$:
\[
  f(\lambda x + (1-\lambda)y) \;\leq\; \lambda f(x) + (1-\lambda)f(y).
\]
Geometrically: the line segment between any two points on the graph lies above the function.
Equivalently, if $f$ is twice differentiable, the Hessian $\nabla^2 f \succeq 0$ everywhere.

\subsection{Q: ``Why is logistic loss convex?''}
The Hessian of logistic loss is $\nabla^2 f = X^\top D X$, where $D$ is diagonal with entries $D_{ii} = \sigma_i(1-\sigma_i)$.
Since $\sigma_i \in (0,1)$, all entries are positive, so $D \succ 0$.
Then $X^\top D X \succeq 0$ for any $X$, which means the loss is convex.

\subsection{Q: ``What makes your non-convex regularizer non-convex?''}
The regularizer is:
\[
  r(w) = \sum_{j} \frac{\alpha\, w_j^2}{1 + \alpha\, w_j^2}.
\]
Its second derivative changes sign: for small $w_j$ it behaves like $\alpha w_j^2$ (convex), but as $w_j$ grows, the term saturates at 1 and the curvature becomes negative.
So the overall objective --- convex logistic loss plus this regularizer --- is non-convex.
We use $\lambda = 0.01$, $\alpha = 1$ following Reddi et al.\ (2016).

\subsection{Q: ``What does $O(1/\sqrt{T})$ convergence mean?''}
The suboptimality $f(w_t) - f^*$ decreases at rate $\propto 1/\sqrt{T}$.
After $T$ iterations, we are at most $c/\sqrt{T}$ away from optimal.
Concretely: to halve the error, you need $4\times$ more iterations.
This is \textbf{sublinear} convergence, typical for stochastic first-order methods.

\subsection{Q: ``Is $O(1/\sqrt{T})$ linear convergence?''}
\textbf{No.} Linear convergence means exponential decrease:
\[
  f(w_t) - f^* \;\leq\; \rho^t \cdot (f(w_0) - f^*), \quad \rho < 1.
\]
On a log-scale plot, linear convergence is a straight line going down.
$O(1/\sqrt{T})$ is sublinear --- it flattens out.
SGD on strongly convex problems \emph{can} achieve linear convergence with decreasing step size, but in the stochastic non-convex setting, $O(1/\sqrt{T})$ is the standard rate.

\subsection{Q: ``Difference between convergence rate and convergence guarantee?''}
The \textbf{rate} tells you \emph{how fast}: e.g.\ $O(1/\sqrt{T})$.
The \textbf{guarantee} specifies \emph{under what conditions}: bounded gradients, Lipschitz smoothness, unbiased gradient estimates.
RMSprop has no formal convergence guarantee --- no proven rate under standard assumptions.

%=============================================================================
\section{Optimizer-Specific Questions}
%=============================================================================

\subsection{Q: ``Difference between Adagrad and AdaGrad-Norm?''}
\textbf{Adagrad:} per-coordinate accumulator $G_{jj} = \sum_t g_{t,j}^2$.
Each coordinate gets its own effective learning rate $\eta / \sqrt{G_{jj}}$.

\textbf{AdaGrad-Norm:} single scalar $G_t = \sum_t \|g_t\|^2$.
All coordinates share the same scaling $\eta / \sqrt{G_t}$.

Adagrad is better when features have heterogeneous scales (sparse data).
AdaGrad-Norm is simpler and has the same $O(1/\sqrt{T})$ rate.

\subsection{Q: ``Why does Adagrad's step size shrink to zero?''}
The accumulator $G_{jj}$ grows monotonically --- it sums squared gradients without forgetting.
After many iterations, $\sqrt{G_{jj}}$ becomes very large and the effective step size $\eta/\sqrt{G_{jj}} \to 0$.
This is actually \emph{desired} for convex problems (you want the step size to decrease to converge), but for non-convex problems or long training runs, it can cause premature convergence.

\subsection{Q: ``Why does Adam sometimes not converge?''}
Reddi et al.\ (2018) showed a concrete counterexample.
The issue: exponential forgetting in the second moment estimate $\hat{v}_t$.
If a rare but large gradient appears, the denominator can become very small, causing the step size to blow up.
Fix: \textbf{AMSGrad} uses $\hat{v}_t = \max(\hat{v}_{t-1}, v_t)$, ensuring the denominator never shrinks.

\subsection{Q: ``Why does RMSprop work so well despite no convergence guarantee?''}
RMSprop uses exponential forgetting --- it only remembers recent gradients, making it very adaptive to changes in the loss landscape.
The lack of formal guarantees means we can't prove a worst-case rate, but empirically it's consistently among the fastest.
That said, it \emph{can} diverge at large learning rates --- we show this in our LR sweep (backup slide~24).

\subsection{Q: ``Is Adagrad always better than SGD?''}
No. Adagrad's advantage comes from per-coordinate scaling, which matters most with sparse or heterogeneously-scaled features.
On dense, uniformly-scaled features, the benefit is small.
Also, for very long training runs, Adagrad's monotonically decreasing step size can be a disadvantage vs.\ SGD with a tuned schedule.

%=============================================================================
\section{Setup \& Methodology}
%=============================================================================

\subsection{Q: ``What model are you using? Is it a neural network?''}
\textbf{No, linear models.}
\begin{itemize}[nosep]
  \item A9a: logistic regression --- linear classifier with sigmoid output.
  \item MNIST: softmax regression --- linear multi-class classifier.
\end{itemize}
No hidden layers, no nonlinear activations.
We chose linear models deliberately: the focus is on comparing \emph{optimizers}, not architectures.
A linear model gives a clean, controlled setting to isolate the optimizer effect.

\subsection{Q: ``Why linear models and not neural networks?''}
Three reasons:
\begin{enumerate}[nosep]
  \item Logistic loss on a linear model is \textbf{convex} $\Rightarrow$ convergence guarantees apply directly.
  \item Isolates the optimizer effect --- no architecture confounds.
  \item Original papers (Duchi et al.\ 2011, Ward et al.\ 2020) also use linear models in their experiments.
\end{enumerate}

\subsection{Q: ``Why these two datasets?''}
\textbf{A9a}: 123 sparse features, 32K samples --- standard binary benchmark where per-coordinate scaling matters due to feature heterogeneity.\\
\textbf{MNIST}: 784 dense features, 60K samples --- higher-dimensional multi-class problem.\\
Together they test adaptive methods in two regimes: sparse vs.\ dense, low vs.\ high dimensionality.

\subsection{Q: ``How did you choose the learning rates?''}
Grid search over $\eta \in \{10^{-4},\, 3\!\times\!10^{-4},\, 10^{-3},\, \ldots,\, 1\}$ --- 9 values on a log scale.
For each optimizer we picked the $\eta$ with the best final training loss.
Backup slide~24 shows the full sweep.
Key finding: Adagrad is very robust across the entire range; RMSprop diverges at $\eta = 1$.

\subsection{Q: ``How many epochs? What batch size?''}
10--15 epochs, mini-batch size 128 (A9a) and 256 (MNIST).
All results averaged over 3 random seeds to reduce noise.

\subsection{Q: ``Did you implement everything from scratch?''}
Yes. All five optimizers implemented from scratch in NumPy.
Loss functions and gradient computations are our own code.
We handle both dense and sparse matrices (\texttt{scipy.sparse} for A9a).

%=============================================================================
\section{Results Interpretation}
%=============================================================================

\subsection{Q: ``Is Adagrad converging faster than the theoretical rate?''}
Yes, significantly.
Theoretical worst-case for Adagrad: $O(\log T / \sqrt{T})$.
Empirically on A9a: the loss drops below $10^{-7}$ suboptimality.
This is expected --- worst-case bounds hold for \emph{arbitrary} convex functions.
Our logistic loss has additional structure (smoothness, bounded spectral norm of $X$).

\subsection{Q: ``Why are the final accuracies so similar across optimizers?''}
All optimizers minimize the \emph{same} convex objective and converge to the \emph{same} global minimum --- just at different speeds.
The final model is essentially identical regardless of optimizer.
\textbf{Key insight:} adaptive methods are about \emph{speed}, not better solutions.

\subsection{Q: ``Why does the non-convex regularizer not change the ranking?''}
At $\lambda = 0.01$, the regularizer is small relative to the logistic loss --- the objective is ``mildly'' non-convex.
In our $\lambda$-sweep (backup slide~26), at $\lambda \geq 0.1$ the effect becomes significant and accuracy drops.
The optimizer ranking stays similar because all methods handle this mild non-convexity similarly.

\subsection{Q: ``Gradient norm plot --- what does it mean when the norm goes to zero?''}
We've reached a stationary point: $\nabla f(w) \approx 0$.
For convex problems $\Rightarrow$ global minimum.
For non-convex $\Rightarrow$ could be saddle point or local minimum.
Key observation: adaptive methods bring the gradient norm down faster --- they normalize the step size, preventing overshooting when gradients are large early on.

%=============================================================================
\section{Theory vs.\ Practice}
%=============================================================================

\subsection{Q: ``Do your empirical results match the theoretical rates?''}
Backup slide~27 shows the comparison.
All methods converge \emph{faster} than worst-case bounds.
Qualitatively, ranking is consistent: Adagrad converges fastest, SGD slowest.
The bounds are for adversarial functions; our problems have nice structure.

\subsection{Q: ``What assumptions do you need for SGD's $O(1/\sqrt{T})$?''}
Standard assumptions:
\begin{enumerate}[nosep]
  \item $L$-Lipschitz continuous gradients: $\|\nabla f(x) - \nabla f(y)\| \leq L\|x - y\|$
  \item Unbiased stochastic gradients: $\mathbb{E}[g_t] = \nabla f(w_t)$
  \item Bounded variance: $\mathbb{E}[\|g_t - \nabla f(w_t)\|^2] \leq \sigma^2$
\end{enumerate}
Under these, SGD with $\eta = c/\sqrt{T}$ achieves $\mathbb{E}[\|\nabla f(w)\|^2] \leq O(1/\sqrt{T})$ for non-convex objectives.

\subsection{Q: ``Why care about convergence rates if all methods reach the same solution?''}
Compute budget is limited.
If method A reaches 90\% accuracy in 100 iterations and B needs 1000, that's $10\times$ more training time.
For large-scale problems, this translates to days vs.\ weeks.
Convergence rates tell us which method uses compute most efficiently.

%=============================================================================
\section{Quick-Fire Definitions}
%=============================================================================

\begin{center}
\renewcommand{\arraystretch}{1.3}
\begin{tabular}{@{}l p{10.5cm}@{}}
\toprule
\textbf{Term} & \textbf{One-line answer} \\
\midrule
Positive definite & All eigenvalues $> 0$; equivalently $x^\top A x > 0$ for all $x \neq 0$ \\
Positive semidefinite & All eigenvalues $\geq 0$; equivalently $x^\top A x \geq 0$ for all $x$ \\
Lipschitz gradient & $\|\nabla f(x) - \nabla f(y)\| \leq L\|x-y\|$ --- gradient doesn't change too fast \\
Unbiased estimator & $\mathbb{E}[g_t] = \nabla f(w_t)$ --- on average, stochastic gradient = true gradient \\
Momentum & Running average of past gradients to smooth updates and accelerate \\
Bias correction (Adam) & Divides by $(1 - \beta^t)$ to compensate for zero-initialization in first steps \\
Sublinear convergence & Error $\to 0$ but slower than geometric. Example: $O(1/\sqrt{T})$ \\
Linear convergence & Error decreases by constant factor each step: $e_t \leq \rho^t e_0$, $\rho < 1$ \\
Superlinear convergence & Faster than linear; ratio $e_{t+1}/e_t \to 0$. Example: Newton's method \\
\bottomrule
\end{tabular}
\end{center}

\end{document}
